\documentclass[12pt]{article}
\usepackage[utf8]{inputenc}
\usepackage[T1]{fontenc}
\usepackage[margin=1in]{geometry}
\usepackage{setspace}
\usepackage{xcolor}
\usepackage[colorlinks=true, urlcolor=blue!60!black, linkcolor=blue!60!black]{hyperref}
\setlength{\parindent}{0em}
\setlength{\parskip}{0.45em}
\renewcommand{\baselinestretch}{1.0}
\begin{document}
\begin{center}
{\Large \textit{Mathews v. Eldridge}, 424 U.S. 319 (1976)\par}
\vspace{0.3em}
{\normalsize Case Brief --- One-Page Edition\par}
\end{center}

\textbf{Memory Jogger.} The three-factor test for procedural due process (\href{https://supreme.justia.com/cases/federal/us/424/319/}{424 U.S. 319}).

\textbf{Facts.} Social Security cut off George Eldridge's disability benefits after a written review of medical reports, notice, and a chance for written rebuttal; an evidentiary hearing and retroactive benefits were available afterward.

\textbf{Procedural History.} Relying on \textit{Goldberg v.\ Kelly}, the lower courts required a pre-termination hearing.

\textbf{Votes.} Reversed, 6--2. Powell wrote for the Court; Brennan dissented, joined by Marshall.

\textbf{Holding.} Due process does not require a hearing before disability benefits end. Courts weigh the private interest, the risk of error and value of added safeguards, and the government's interest.

\textbf{Reasoning.} Unlike welfare in \textit{Goldberg}, disability benefits are not need-based and back pay is available; medical determinations rest on standard reports rather than credibility, so written procedures limit error; and extra hearings would drain limited funds. Brennan stressed the real hardship to Eldridge's family.

\textbf{Examples.} \emph{Same-Side:} reducing a veteran's rating on medical evidence. \emph{Opposite-Side:} cutting off welfare over a credibility dispute. \emph{Fence-Sitter:} algorithmic Medicaid cuts.

\textbf{Critique.} Mashaw says the test ignores dignitary values; originalists call it ahistorical.

\textbf{Key Quotation.} ``More precisely, our prior decisions indicate that identification of the specific dictates of due process generally requires consideration of three distinct factors: first, the private interest that will be affected by the official action; second, the risk of an erroneous deprivation of such interest through the procedures used, and the probable value, if any, of additional or substitute procedural safeguards; and, finally, the Government's interest, including the function involved and the fiscal and administrative burdens that the additional or substitute procedural requirement would entail.'' --- Justice Powell, Opinion of the Court, 424 U.S. at 334--35.

\end{document}
